\section{Experimental Setup}

\subsection{Evaluation Protocol}

The experiments address seven questions: task learnability (RQ1), human benchmark performance (RQ2), component contributions (RQ3), the effect of peptide-disjoint evaluation (RQ4), replication in the mouse tissue--H2 benchmark (RQ5), how much signal remains without tissue identity (RQ6), and whether architecture advantages persist under the strict peptide-disjoint protocol (RQ7). Cross-species comparisons concern task-level replication, not parameter transfer.

The project includes extensive model development on the standard human benchmark. To reduce further selection bias, the strict analyses are governed by a frozen confirmatory protocol. Before inspecting peptide-disjoint performance, the protocol fixes the primary and secondary metrics, split construction, split seed, model seeds, task-inclusion rules, model configurations, training budgets, fusion and ensemble rules, baseline set, and statistical comparisons. Split manifests, configuration hashes, and timestamps are retained with the resulting predictions.

The primary standard results and strict robustness results answer different questions. The existing human standard pair-disjoint result is the primary closed-task benchmark result, but it reflects a long model-development history and is reported with that limitation. The mouse standard result uses a model, seed set, and probability-averaging rule frozen before a one-time internal fixed-test evaluation. The connected-component peptide-disjoint out-of-fold experiments in both species are confirmatory evidence about seen-task, unseen-peptide behavior. They supplement rather than replace the standard results.

For every trainable method, each held-out fold is isolated from model fitting. Encoders, heads, preprocessing state, auxiliary classifiers, early-stopping rules when applicable, and any data-dependent thresholds are estimated exclusively from the fitting partition. The final held-out prediction for an example is generated only by models trained without its connected peptide component. Model selection is never performed on pooled strict out-of-fold predictions.

A matched standard pair-disjoint control is constructed from the same pair pool used by the strict protocol. The number of folds, task coverage, and task-wise held-out counts are matched to the peptide-disjoint folds as closely as possible. Only this matched comparison is used to attribute a performance change to peptide separation. Differences between the earlier fixed tests and the strict out-of-fold experiments are reported descriptively because the underlying sample pools and split constructions differ.

\subsection{Training and Baselines}

All encoders and heads are optimized with AdamW \cite{loshchilov2019adamw} for 25 epochs using a batch size of 512, a learning rate of \(10^{-3}\), and weight decay of \(10^{-4}\). Gradient norms are clipped at 1.0. The global and HLA-specific branches are trained independently for each random seed. Within every out-of-fold experiment, only the fitting partition is used to estimate model parameters; held-out rows are used solely for prediction and evaluation.

The human main model uses seeds 20260704, 20260705, and 20260706. The architecture, optimizer settings, auxiliary weights, fusion rule, and seed set are fixed for the reported evaluation. Because the training schedule uses a fixed number of epochs, the held-out fold is not used for early stopping. The same preprocessing and task mappings are reconstructed from each fitting partition so that no held-out labels influence training.

The reported human standard comparisons comprise one-hot logistic regression, a shared encoder with task-specific heads, the auxiliary dual branch, a three-seed MLP dual-branch ensemble, and the TissuePMHC human model. The strict architecture comparison expands this set to one-hot logistic regression, a BLOSUM62 random forest, shared heads, a plain MLP dual branch, an auxiliary MLP dual branch, and the TissuePMHC human model. The mouse comparisons comprise a BLOSUM62 random forest, shared heads, Factorized MMoE, and, for the standard protocol only, an H2-Kk residual adapter. All trainable methods use the same split assignments within a comparison and the recorded training seed sets. The standard benchmark tables report only completed runs; methods without held-out predictions are not listed as experimental evidence.

The BLOSUM62 random-forest baseline \cite{henikoff1992blosum} flattens the 20 substitution scores at each of the nine peptide positions into 180 features and fits a separate model per task, using 200 trees, square-root feature subsampling, and a minimum leaf size of two. The MLP dual-branch baselines use 16-dimensional residue embeddings followed by flattening and two 128-unit ReLU layers, with dropout 0.2 after the first layer. The auxiliary version applies tissue and HLA auxiliary losses weighted 0.1 each in its global branch; the plain version removes those auxiliary objectives while retaining the branch structure and fixed equal-weight task-rank fusion. Human neural ensembles use seeds 20260704--20260706. Mouse strict ensembles use seeds 20260704--20260708. Historical optimization baselines use FAMO \cite{liu2023famo}, conflict-averse gradient descent (CAGrad) \cite{liu2021cagrad}, and dual-balancing multi-task learning (DB-MTL) \cite{lin2023dbmtl}; HLA pseudo-sequence conditioning follows the representation principle introduced by NetMHCpan \cite{nielsen2007netmhcpan}. The historical pair-ranking control replaces row-wise classification with a loss that favors the positive member of each matched pair.

Every strict architecture reads the existing component manifest rather than regenerating folds. The human comparison contains 73,798 pairs in 17,036 connected components across 157 tasks; the mouse comparison contains 6,766 pairs in 1,844 components across 24 tasks. Both use three outer folds with zero pair and exact-peptide overlap between fitting and held-out partitions. This matched design separates architecture comparisons from split variation.

\subsection{No-tissue and Frozen General-pMHC Controls}

The strongest direct control for tissue identity is a capacity-matched model that receives peptide and MHC information but no tissue input. For human, this completed control uses the same multi-kernel encoder family as the human main run, has 137,443 trainable parameters, and averages three fixed seeds. It removes tissue identity, tissue embeddings, tissue auxiliary loss, and the tissue--MHC task head, so its score has the form \(s=f(p,m)\). For mouse, the completed control retains the same MLP and three-expert configuration as the mouse main run but gates only on peptide and MHC, has 65,943 trainable parameters, and averages five fixed seeds. These models use the same standard and peptide-component folds as the full models; the reported analyses use their completed row-level prediction archives.

To test whether performance is primarily explained by generic peptide--MHC propensity, we score every unique peptide--MHC query with two external models that do not read tissue identity and are not fine-tuned on the benchmark labels. MHCflurry 2.2.1 \cite{odonnell2020mhcflurry} is evaluated using its presentation score with peptide-flank features disabled. NetMHCpan 4.1b \cite{reynisson2020netmhcpan} is evaluated primarily using eluted-ligand rank (EL rank), with binding-affinity rank (BA rank) retained as a binding-only sensitivity analysis. Rank outputs are sign-reversed so that larger values consistently indicate stronger predicted propensity.

The control set contains 79,759 unique human peptide--MHC queries across 35 HLA alleles and 6,663 unique mouse queries across four H2 restrictions. MHCflurry presentation, NetMHCpan EL, and NetMHCpan BA all achieve 100\% row, pair, and task coverage, so comparisons use exactly the same evaluable pairs as the tissue-conditioned models. Controls are evaluated on the standard fixed tests and on the complete train pools used for matched standard and connected-component peptide-disjoint OOF. Because the external scores are frozen and both OOF protocols use the same train-row pool, their standalone OOF metrics are identical across those two protocols; the protocol-specific comparison is against the corresponding tissue-conditioned OOF predictions.

Incremental information is assessed with a cross-fitted logistic stacker. The external-only stack uses one frozen external score, whereas the combined stack uses that score together with the OOF or fixed-test score of the tissue-conditioned model. Stacker coefficients and any preprocessing are estimated without the rows on which the stacked prediction is evaluated. This analysis tests residual information in the external score; it is not treated as a newly optimized main architecture.

\subsection{Metrics and Statistical Analysis}

The primary metric is task-macro AUROC:
\[
\operatorname{MacroAUROC}
=
\frac{1}{|\mathcal{T}_{\mathrm{valid}}|}
\sum_{\tau\in\mathcal{T}_{\mathrm{valid}}}
\operatorname{AUROC}_{\tau},
\]
where \(\mathcal{T}_{\mathrm{valid}}\) contains the prespecified tasks with both labels available for evaluation. Macro averaging gives every task equal weight regardless of its sample size. Task-macro AUPRC is defined analogously:
\[
\operatorname{MacroAUPRC}
=
\frac{1}{|\mathcal{T}_{\mathrm{valid}}|}
\sum_{\tau\in\mathcal{T}_{\mathrm{valid}}}
\operatorname{AUPRC}_{\tau}.
\]
Because the paired dataset has a constructed 1:1 label ratio, AUPRC is interpreted relative to the per-task benchmark prevalence and is not compared directly with AUPRC from naturally imbalanced clinical cohorts.

To measure lower-tail performance, we report the mean AUROC of the ten lowest-performing human tasks and the six lowest-performing mouse tasks. PairAcc, defined in the benchmark section, measures strict concordance within the original matched pairs: a score tie receives zero credit. As a sensitivity analysis, assigning half credit to ties changes the human matched-standard and strict values by only 0.0003 and 0.0002, respectively, and does not change mouse values because no tied pairs occur. For every Accuracy, balanced-accuracy, F1, and MCC value in the reported benchmark tables, the saved score is converted to a label by the fixed rule \(\hat y=\mathbf{1}[s\geq0.5]\); this threshold is not optimized on fitting, validation, or held-out labels. For rank-fused models, 0.5 is the midpoint of a within-batch percentile-rank score rather than a calibrated probability cutoff. AUROC, AUPRC, and PairAcc remain the threshold-free ranking metrics used for the primary conclusions.

Metrics are computed in the following order. For each seed, predictions from all outer held-out folds are first concatenated. Each task metric is then computed from its complete out-of-fold prediction vector, after which task metrics are averaged with equal weight. Strict ensemble predictions for a row are produced only from members trained on the corresponding outer fitting fold. Fold-level metrics are reported descriptively but are not averaged as if folds were independent experimental replicates. Single-seed means and standard deviations quantify training variation, whereas the ensemble remains a single final predictor.

Strict-protocol uncertainty is estimated by 10,000 task-bootstrap replicates with seed 20260711. The bootstrap resamples tasks with replacement and therefore characterizes task-level uncertainty without treating outer folds or ensemble members as independent replicates. Standard-versus-strict, strict architecture, and tissue-conditioned-versus-external control differences are evaluated with two-sided Wilcoxon signed-rank tests; we report the median and Hodges--Lehmann differences, win/tie/loss counts, percentile confidence intervals, and Benjamini--Hochberg adjusted values within their prespecified species--metric analysis families. Pearson correlation is used for within-task branch-score agreement, whereas Spearman correlation is used for associations involving task performance, training size, or fusion gain. Subgroup and motif-related observations are treated as descriptive unless a prespecified inferential test is reported. Tasks are the analysis units but are not fully independent because multiple tasks can share a tissue, an MHC restriction, parent proteins, and peptide connected components. The task bootstrap and Wilcoxon test do not model these crossed dependencies or model-retraining variance. Their intervals and adjusted \(p\) values are therefore nominal task-level summaries and may be optimistic; we emphasize effect sizes, direction counts, and consistency across species rather than treating them as independent-cohort inference.

\subsection{Cross-species Evaluation and Reproducibility}

The human and mouse benchmarks share the paired task definition, row labels, split principles, primary metric, and leakage audits. They differ in scale, MHC system, and selected model. Human experiments center on the TissuePMHC human model across 157 tissue--HLA tasks. Mouse experiments center on a lightweight Factorized MMoE across 24 tissue--H2 tasks, followed by equal-probability averaging across five frozen seeds.

The human model uses seeds 20260704, 20260705, and 20260706. The mouse ensemble uses those three seeds plus 20260707 and 20260708. Pair-grouped and peptide-component fold assignments use the recorded split seed 20260711. Comparisons across species emphasize whether both benchmarks retain above-random signal and whether structured sharing improves upon simple sharing under matched protocols. We do not compare the numerical advantage of the two architectures as a controlled species effect, because data size, task diversity, and model selection differ.

For each experiment, we retain the exact train/test access policy, split manifest, fold seed, training seeds, model configuration, and command-line arguments. Reports include epoch count, batch size, optimizer, learning rate, weight decay, dropout, parameter count, training time, peak memory when available, hardware, Python and PyTorch versions, and random-seed handling. Predictions are saved at the individual-row level with task, label, pair identifier, fold, model member, and final score whenever supported by the run.

Dataset and configuration files are accompanied by hashes, and each result directory records the timestamped protocol and code state. The standard mouse fixed-test metadata explicitly records that test data were accessed only for the frozen evaluation and that model selection was not performed on the test set. Peptide-disjoint runs retain their component assignments and zero-overlap audits.

\subsection{Data and Code Availability}

The accompanying project repository contains the processed benchmark tables, pair and component split manifests, row-level predictions used for the reported analyses, configuration files, hashes, and analysis scripts. Raw IEDB-derived material and third-party predictor assets remain subject to their original access and licensing terms; NetMHCpan executables are not redistributed.
